%---------------------------------------------------------------
\section{Nasazení aplikace}\label{sec:deployment}
%---------------------------------------------------------------

% Review
V rámci nasazování aplikace jsem vytvořil tři samostatná prostředí \texttt{dev}, \texttt{stage} a \texttt{production}. Každé prostředí má odlišný účel, odlišné napojení na API a odlišný způsob nasazení. Kromě nasazení aplikace bylo třeba také zajistit, aby byly k dispozici pro stažení a instalaci vytvořené knihovny \texttt{eduriam/ui-core} a \texttt{eduriam/ui-x}. Vše podrobněji popíšu v této podkapitole.

% Review
\subsection{Vývojové prostředí (\texttt{dev})}

% Review
Testovací prostředí \texttt{dev} je dostupné na adrese \texttt{dev.eduriam.com}. Hosting je zajištěn na platformě Vercel. Tento hosting jsem zvolil jednak protože původní aplikace tento hosting využívala a jednak protože společnost Vercel vyvíjí framework Next.js, ve kterém je tato aplikace naprogramována. Aplikace v tomto prostředí využívá mock backend API, který je nasazený na platformě Netlify.

% Review
Nasazení do tohoto prostředí probíhá automaticky při každém provedení \texttt{push} do větve \texttt{dev} v repozitáři frontendu aplikace. Při provedení této akce se provede jak nasazení aplikace, tak nasazení zmíněného mocku API.

% Review
Toto prostředí slouží k několika účelům. Jednak slouží k rychlému vyvíjení a testování funkcí aplikace, které ještě nejsou podporovány backend API. Dále lze toto prostředí snadno využít pro uživatelské testování a validaci nových funkcí.

% Review
\subsection{Testovací prostředí (\texttt{stage})}

% Review
Testovací prostředí \texttt{stage} je dostupné na adrese \texttt{app.stage.eduriam.com} a rovněž běží na hostingu Vercel. Na rozdíl od prostředí \texttt{dev} je toto prostředí napojeno na stage backendové API, tedy na reálnou implementaci API určenou pro testování.

% Review
Nasazení frontendu do stage probíhá také automaticky při každém provedení \texttt{push} do větve \texttt{dev}. Nasazením backendu se zabývá kolega Bc. Jakub Vondráček v jeho diplomové práci.

% Review
Stage prostředí tak představuje hlavní testovací prostředí, ve kterém lze ověřovat chování aplikace se skutečnými koncovými body API.

% Review
\subsection{Produkční prostředí (\texttt{production})}

% Review
Produkční prostředí \texttt{production} je dostupné na adrese \texttt{app.eduriam.com} a je napojeno na produkční backendové API. Podobně jako ostatní prostředí je toto prostředí nasazeno na platformě Vercel.

% Review
Nasazení probíhá při manuálním vytvoření release značky na větvi \texttt{master}, což odpovídá procesu nasazování, který byl zaveden v původní aplikaci. Nasazením produkčního API se opět zabývá kolega Bc. Jakub Vondráček v jeho diplomové práci.

% Review
Produkční prostředí tak představuje hlavní prostředí, ve kterém je aplikace veřejně dostupná uživatelům.

% Review
\subsection{Nasazení knihoven aplikace \texttt{@eduriam/ui-core} a \texttt{@eduriam/ui-x}}

% Review
Kromě nasazení frontendu bylo třeba zajistit, aby byly k dispozici pro stažení a instalaci nově vytvořené knihovny \texttt{@eduriam/ui-core} a \texttt{@eduriam/ui-x}. Knihovny jsem se rozhodl hostovat na platformě GitHub s využitím služby GitHub Packages\footnote{Dostupné na: \href{https://docs.github.com/packages}{https://docs.github.com/packages}}. Tato platforma umožňuje snadno hostovat a distribuovat knihovny sestavením balíčků přímo z GitHub repozitáře. 

% Review
Nasazení nové verze probíhá automaticky při vytvoření \texttt{merge} z větve \texttt{dev} do větve \texttt{main} v repozitáři knihoven, pokud byla v rámci změn navýšena verze knihovny v souboru \texttt{package.json}. Po nasazení lze knihovny stáhnout a instalovat pomocí nástrojů \texttt{npm} nebo \texttt{yarn} podobně jako ostatní knihovny.
